\documentclass[11pt]{article}

% Packages
\usepackage[utf8]{inputenc}
\usepackage[T1]{fontenc}
\usepackage{amsmath,amssymb,amsthm}
\usepackage{mathtools}
\usepackage{hyperref}
\usepackage{cleveref}
\usepackage{booktabs}
\usepackage{enumitem}
\usepackage[margin=1in]{geometry}
\usepackage{xcolor}

% Theorem environments
\newtheorem{theorem}{Theorem}[section]
\newtheorem{lemma}[theorem]{Lemma}
\newtheorem{corollary}[theorem]{Corollary}
\newtheorem{proposition}[theorem]{Proposition}
\newtheorem{definition}[theorem]{Definition}
\newtheorem{remark}[theorem]{Remark}
\newtheorem{example}[theorem]{Example}
\newtheorem{counterexample}[theorem]{Counterexample}

% Commands
\newcommand{\R}{\mathbb{R}}
\newcommand{\N}{\mathbb{N}}
\newcommand{\cl}[1]{\overline{#1}}
\newcommand{\Fix}{\operatorname{Fix}}
\newcommand{\AD}{\operatorname{AD}}
\newcommand{\id}{\operatorname{id}}
\newcommand{\diam}{\operatorname{diam}}

\title{The Defense Impossibility Theorem:\\
  Continuous Utility-Preserving Defenses Fix Boundary Points}

\author{
  Manish Bhatt\thanks{Corresponding author: \texttt{manish.bhatt13212@gmail.com}. This work was conducted independently and does not reflect the views of the authors' respective employers.} \\
  \and
  Sarthak Munshi \\
  \and
  Vineeth Sai Narajala \\
  \and
  Idan Habler \\
  \and
  Ammar Al-Kahfah \\
  \and
  Ken Huang \\
  \and
  Blake Gatto \\
}

\date{}

\begin{document}

\maketitle

\begin{abstract}
We prove that no continuous, utility-preserving defense can fully eliminate
prompt injection vulnerabilities at the safety boundary of a language model.
Specifically, if the alignment deviation function $f\colon X \to \R$ is
continuous on a connected Hausdorff space~$X$ and a defense
$D\colon X \to X$ is continuous with $D(x) = x$ for all safe inputs, then
there exist boundary points~$z$ with $f(z) = \tau$ that $D$ is forced to
fix in place. The proof requires five steps and no machinery beyond
point-set topology; it is machine-verified in Lean~4 with Mathlib
(zero \texttt{sorry}, three standard axioms). Beyond the central
impossibility, we derive nine supporting results that characterize the
vulnerability landscape---including basin structure,
perturbation robustness, iterative convergence, cross-model transferability,
authority monotonicity, gradient ascent, and exponential cost
asymmetry---forming a complete formal theory of the Manifold of Failure.
We connect each theorem to empirical findings from behavioral mapping of
three frontier LLMs (Llama-3-8B, GPT-OSS-20B, GPT-5-Mini), and find
theory--experiment agreement at every point of contact. The complete Lean
artifact (33 files, $\sim$200 theorems) is open-sourced.
\end{abstract}

% ============================================================
\section{Introduction}
\label{sec:intro}

Can you build a wrapper around a language model that eliminates all
prompt injection vulnerabilities?

Most current work on LLM safety implicitly assumes the answer is
yes---at least in principle, even if current tools fall short in practice.
The defense mechanisms people have proposed span a wide range: input
classifiers that flag suspicious prompts~\cite{alon2023detecting},
output filters that catch unsafe completions~\cite{inan2023llama},
and constitutional rewriting pipelines that reshape model
behavior~\cite{bai2022constitutional}. But despite their differences,
all of these defenses share the same underlying structure: they are
functions $D\colon X \to X$ that preprocess prompts before the model
sees them, aiming to map unsafe inputs to safe equivalents while leaving
safe inputs alone.

We show that this goal is mathematically impossible under three mild
conditions:
\begin{enumerate}[nosep]
  \item The defense is \emph{continuous}: small changes to the input
    produce small changes to the defense's output.
  \item The defense is \emph{utility-preserving}: it does not modify
    inputs that are already safe.
  \item The prompt space is \emph{connected}: you can continuously
    interpolate between any two prompts.
\end{enumerate}

\noindent When all three hold, the defense is forced to leave at least one
point on the safety boundary completely untouched---a prompt where
alignment deviation sits exactly at the threshold, passed through
with no remediation applied.

The proof itself is short (five steps), uses only elementary point-set
topology, and has been machine-verified in Lean~4 with zero
\texttt{sorry} statements. All three conditions are individually
necessary: we give counterexamples showing the theorem breaks down
whenever any single one is dropped.

It is worth being clear about what this result does \emph{not} say.
It does not say that defense is hopeless. What it says is that
\emph{perfect} defense---a wrapper that renders every possible input
safe---is impossible when the defense is continuous and utility-preserving.
And the impossibility lives at the boundary, which is the mildest possible
kind of failure. The practical takeaway is not that we should give up on
the boundary, but that we should make it shallow: set the threshold
high enough that boundary-level behavior is benign rather than dangerous.
We show that GPT-5-Mini's empirically observed hard ceiling at
$\AD = 0.50$~\cite{bhatt2025manifold} is exactly this strategy at work.

Around this central impossibility result, we develop a broader formal
theory of the Manifold of Failure. Nine supporting theorems characterize
the vulnerability landscape from multiple angles: its topology, how
robust it is to perturbation, why iterative attacks converge, when
attacks transfer across models, how authority framing affects
vulnerability, how gradient-based methods navigate the space, and
why defense costs grow exponentially faster than attack costs.
Each theorem is grounded in empirical findings from behavioral mapping
of three frontier LLMs.

\paragraph{Contributions.}
\begin{itemize}[nosep]
  \item A formally verified impossibility theorem for continuous
    utility-preserving defenses (\Cref{thm:main}).
  \item A five-step proof relying only on Hausdorff separation,
    connectedness, and continuity, together with counterexamples showing
    each hypothesis is necessary (\Cref{sec:proof}).
  \item Nine supporting theorems that collectively form a formal theory of
    the failure manifold (\Cref{sec:landscape,sec:attacks,sec:cost}).
  \item Theory--experiment correspondence with empirical findings from
    three LLMs (\Cref{sec:experiments}).
  \item A Lean~4 artifact comprising 33 files, $\sim$200 theorems, zero
    \texttt{sorry}, and three axioms (\Cref{sec:artifact}).
\end{itemize}


% ============================================================
\section{Related Work}
\label{sec:related}

\paragraph{Adversarial robustness.}
The existence of adversarial examples in neural networks has been known
for over a decade~\cite{szegedy2013intriguing,goodfellow2014explaining},
and a substantial body of work has pursued formal verification of
robustness properties---through abstract
interpretation~\cite{singh2019abstract}, SMT-based
methods~\cite{katz2017reluplex}, and randomized
smoothing~\cite{cohen2019certified}. These approaches verify that
particular inputs are robust to perturbation. Our work operates at a
different level: rather than checking individual inputs, we prove a
structural impossibility for the defense mechanism itself.

\paragraph{Prompt injection and jailbreaking.}
A variety of techniques have been developed to find individual
vulnerabilities in language models, including gradient-based attacks
like GCG~\cite{zou2023universal}, search-based approaches like
PAIR~\cite{chao2024jailbreaking} and TAP~\cite{mehrotra2024tree},
and quality-diversity methods like Rainbow
Teaming~\cite{samvelyan2024rainbow}. What these methods have in common
is that they discover specific failure points. The Manifold of Failure
framework~\cite{bhatt2025manifold} takes a different tack, using
MAP-Elites to map the \emph{continuous topology} of failures rather
than cataloguing individual instances. Our work provides the formal
theory that underpins that framework.

\paragraph{Impossibility results in ML.}
The no-free-lunch theorems~\cite{wolpert1997no} set fundamental limits
on what optimization can achieve in general, and the robustness--accuracy
tradeoff literature~\cite{tsipras2018robustness} demonstrates that
pursuing adversarial robustness comes at a cost to standard performance.
Our result is different in character from both: it is a topological
impossibility, not a statistical tradeoff. It holds for \emph{any}
continuous function on \emph{any} connected Hausdorff space, with no
dependence on the learning algorithm, training data, or distribution.

\paragraph{Formal verification of ML systems.}
Proof assistants like Lean and Coq have been applied to verify properties
of neural networks~\cite{bagnall2019certifying} and optimization
algorithms~\cite{bolte2024}. To the best of our knowledge, ours is the
first formally verified impossibility result specifically targeting
prompt injection defense.


% ============================================================
\section{Setup}
\label{sec:setup}

We begin by establishing the formal objects we will work with throughout
the paper.

\begin{definition}[Alignment Deviation Function]
\label{def:ad}
Let $X$ be a topological space. An \emph{alignment deviation function} is
a continuous map $f\colon X \to \R$. Given a threshold $\tau\in\R$, we
partition the space into three regions:
\begin{align}
  S_\tau &= \{x \in X : f(x) < \tau\}  &&\text{(safe region)} \\
  U_\tau &= \{x \in X : f(x) > \tau\}  &&\text{(unsafe region)} \\
  B_\tau &= \{x \in X : f(x) = \tau\}  &&\text{(boundary)} \\
  F_\tau &= \{x \in X : f(x) \geq \tau\} &&\text{(failure manifold)}
\end{align}
\end{definition}

\begin{definition}[Defense]
\label{def:defense}
A \emph{defense} is a continuous map $D\colon X \to X$. We call it
\emph{utility-preserving} if it acts as the identity on safe inputs:
$D(x) = x$ for all $x \in S_\tau$.
\end{definition}

\begin{definition}[Complete Defense]
\label{def:complete}
A defense is \emph{complete} if $f(D(x)) < \tau$ for all $x\in X$---that
is, after preprocessing, every input lands in the safe region.
\end{definition}


% ============================================================
\section{The Defense Impossibility Theorem}
\label{sec:proof}

\begin{theorem}[Boundary Fixation]
\label{thm:main}
Let $X$ be a connected Hausdorff space. Let $f\colon X\to\R$ be
continuous with $S_\tau \neq \emptyset$ and $U_\tau \neq \emptyset$. Let
$D\colon X \to X$ be continuous with $D|_{S_\tau} = \id$.

Then there exists $z\in X$ with $f(z) = \tau$ and $D(z) = z$.

Moreover, every $z \in \cl{S_\tau} \setminus S_\tau$ satisfies
$f(z) = \tau$ and $D(z) = z$. This set is nonempty.
\end{theorem}

\begin{proof}
The proof proceeds in five steps, each using exactly one of the
hypotheses.

\textbf{Step~1} (The fixed-point set is closed).
Let $\Fix(D) = \{x \in X : D(x) = x\}$. Consider the map
$\varphi\colon X \to X \times X$ defined by $\varphi(x) = (D(x), x)$,
which is continuous since $D$ is. The diagonal
$\Delta = \{(y,y) : y \in X\}$ is closed in $X\times X$ because $X$ is
Hausdorff---this is precisely where we use the Hausdorff hypothesis.
Since $\Fix(D) = \varphi^{-1}(\Delta)$ is the preimage of a closed set
under a continuous map, it is closed.

\textbf{Step~2} (The closure of the safe region consists of fixed points).
Utility preservation tells us that $S_\tau \subseteq \Fix(D)$. Since
$\Fix(D)$ is closed by Step~1, taking closures preserves the inclusion:
$\cl{S_\tau} \subseteq \Fix(D)$.

\textbf{Step~3} (The safe region is not closed).
Because $f$ is continuous, $S_\tau = f^{-1}((-\infty,\tau))$ is the
preimage of an open set, hence open. Now suppose for contradiction that
$S_\tau$ were also closed. Then it would be clopen, and in a connected
space the only clopen sets are $\emptyset$ and $X$ itself. But
$S_\tau \neq \emptyset$ by hypothesis, and $S_\tau \neq X$ since
$U_\tau \neq \emptyset$. This is a contradiction, so $S_\tau$ cannot
be closed.

\textbf{Step~4} (A boundary point lives in the closure).
Since $S_\tau$ is not closed (Step~3), its closure strictly contains it:
$\cl{S_\tau} \supsetneq S_\tau$. Pick any
$z \in \cl{S_\tau} \setminus S_\tau$. We claim $f(z) = \tau$.
\begin{itemize}[nosep]
  \item On one hand, $f(z) \leq \tau$: the sublevel set
    $\{x : f(x) \leq \tau\}$ is closed (as the preimage of
    $(-\infty,\tau]$) and contains $S_\tau$, so by minimality of closure
    it contains $\cl{S_\tau}$.
  \item On the other hand, $f(z) \geq \tau$: since
    $z \notin S_\tau$, we know $f(z) \not< \tau$.
\end{itemize}
Together these give $f(z) = \tau$, so $z$ lies on the boundary.

\textbf{Step~5} (The defense fixes this boundary point).
From Step~2, $z \in \cl{S_\tau} \subseteq \Fix(D)$, so $D(z) = z$.
Combined with $f(z) = \tau$ from Step~4, we get
$f(D(z)) = f(z) = \tau$---the defense leaves this boundary point
exactly where it found it.
\end{proof}

\begin{corollary}
\label{cor:no-complete}
Under the hypotheses of \Cref{thm:main}, no continuous utility-preserving
defense is complete.
\end{corollary}

\begin{proof}
Completeness requires $f(D(z)) < \tau$ for every~$z$. But
\Cref{thm:main} gives us a $z$ with $f(D(z)) = \tau \not< \tau$.
\end{proof}


% ============================================================
\subsection{Each Hypothesis Is Necessary}
\label{sec:counterexamples}

To confirm that the theorem is tight, we show that dropping any single
hypothesis allows a complete defense to exist.

\begin{counterexample}[Removing connectedness]
\label{ce:connected}
Take $X = \{0,1\}$ with the discrete topology, so $X$ is disconnected.
Set $f(0)=0$, $f(1)=1$, and $\tau=0.5$. The defense $D(0)=0$, $D(1)=0$
is continuous (every function on a discrete space is), utility-preserving
($D(0)=0$ and $f(0)=0<0.5$), and complete ($f(D(1))=f(0)=0<0.5$). It
works precisely because the safe and unsafe inputs live on disconnected
components with no boundary between them.
\end{counterexample}

\begin{counterexample}[Removing continuity of $D$]
\label{ce:continuous}
Take $X=[0,1]$, $f(x)=x$, and $\tau=0.5$. Define $D(x) = x$ for
$x<0.5$ and $D(x) = 0$ for $x \geq 0.5$. This defense is
utility-preserving and complete ($f(D(x)) < 0.5$ for all~$x$), but it
has a jump discontinuity at $x=0.5$. The defense succeeds because it is
allowed to make an abrupt jump at the boundary.
\end{counterexample}

\begin{counterexample}[Removing utility preservation]
\label{ce:utility}
Fix a safe point $x_0$ and define $D(x) = x_0$ for all~$x$. This
constant map is continuous and complete ($f(D(x)) = f(x_0) < \tau$), but
it modifies every input except $x_0$ itself, destroying utility entirely.
The defense works because it is free to mangle safe inputs.
\end{counterexample}

\begin{remark}[The defense trilemma]
\label{rem:trilemma}
These counterexamples, together with the main theorem, reveal a trilemma.
A defense can satisfy at most two of the following three properties:
\begin{enumerate}[nosep]
  \item \emph{Continuity}: small input changes $\Rightarrow$ small output changes.
  \item \emph{Utility preservation}: safe inputs pass through unchanged.
  \item \emph{Completeness}: all outputs are safe after preprocessing.
\end{enumerate}
\Cref{thm:main} establishes this when the space is connected and both
safe and unsafe inputs exist. Each counterexample achieves two properties
by sacrificing the third.
\end{remark}


% ============================================================
\subsection{What the Theorem Does Not Say}
\label{sec:caveats}

It is equally important to understand the limits of the result.

\begin{remark}[Not all boundary points are fixed]
\label{rem:isolated}
The theorem applies specifically to $\cl{S_\tau} \setminus S_\tau$, not
to the entire boundary $B_\tau$. Boundary points that are ``isolated''
from the safe side---where $f(z)=\tau$ but $f>\tau$ throughout a
punctured neighborhood---need not lie in $\cl{S_\tau}$ and may escape
fixation. The Lean artifact includes a verified example illustrating this:
$f(x)=x^2$ with $\tau=0$, where $f(0)=0$ but
$\{f<0\}=\emptyset$, so $\cl{S_0}=\emptyset$ and the origin is not
captured by the theorem.
\end{remark}

\begin{remark}[Boundary behavior may be mild]
\label{rem:mild}
The fixed points satisfy $f(z) = \tau$ exactly---they sit at the
threshold, not deep in dangerous territory. If the threshold is set high
enough that $\tau$-level behavior is essentially benign, then the theorem
is mathematically true but practically harmless. This is not a loophole
in the result; it \emph{is} the engineering prescription.
\end{remark}

\begin{remark}[Discontinuous defenses escape the theorem]
\label{rem:discontinuous}
Hard blocklists, keyword filters, and exact-match rejection are all
discontinuous operations, and the theorem says nothing about them. Of
course, discontinuous defenses come with their own costs---adversarial
examples crafted to evade the filter, and false positives on benign
inputs---but those are different problems from the one studied here.
\end{remark}


% ============================================================
\section{The Vulnerability Landscape}
\label{sec:landscape}

The Boundary Fixation theorem tells us that the defender cannot clear the
boundary. The natural next question is: what does the landscape around
that boundary actually look like? The following theorems provide the
answer.

\begin{theorem}[Topological Partition]
\label{thm:partition}
For any continuous $f\colon X\to\R$: $S_\tau$ is open, $U_\tau$ is open,
$B_\tau$ is closed, $F_\tau$ is closed, and
$F_\tau = B_\tau \cup U_\tau$. Every point belongs to exactly one of
$S_\tau$, $B_\tau$, $U_\tau$.
\end{theorem}

\begin{theorem}[Basin Structure]
\label{thm:basin}
If $f(p)>\tau$ for some~$p$, then $U_\tau$ is an open set with positive
Lebesgue measure.
\end{theorem}

The significance of this result is that vulnerabilities are never isolated
points. Whenever a single attack succeeds, an entire neighborhood of
attack variants succeeds as well, forming a basin of positive volume.

\begin{theorem}[No Gap]
\label{thm:nogap}
If $X$ is connected and $f$ takes values both above and below~$\tau$,
then $B_\tau \neq \emptyset$. Any continuous path from a safe point to an
unsafe point must cross~$B_\tau$.
\end{theorem}

In other words, there is no moat separating safe from unsafe.
The boundary is a concrete, locatable surface that any continuous
trajectory between the two regions is forced to pass through.


% ============================================================
\section{Attack Properties}
\label{sec:attacks}

With the landscape established, we turn to what it implies about attacks.

\begin{theorem}[Perturbation Robustness]
\label{thm:lipschitz}
If $f$ is $L$-Lipschitz and $f(p)>\tau$, then
$\operatorname{ball}(p,\, r(p)) \subseteq U_\tau$ where
\begin{equation}
  r(p) = \frac{f(p) - \tau}{L} > 0.
\end{equation}
Moreover, $r$ is monotonically increasing in~$f(p)$.
\end{theorem}

This means that successful attacks are robust to rephrasing: there is a
ball of quantifiable radius around any working attack within which every
variant also works. The more severe the attack, the larger this
perturbation ball.

\begin{theorem}[Iterative Convergence]
\label{thm:convergence}
Let $T\colon X\to X$ satisfy
$\operatorname{score}(T(x)) \geq \operatorname{score}(x)$ with
$\operatorname{score}\colon X\to[0,1]$. Then
$\operatorname{score}(T^n(x_0))$ converges. If each step improves the
score by at least $\delta > 0$, convergence occurs in at most
$\lceil 1/\delta \rceil$ steps.
\end{theorem}

This explains a pattern observed across very different attack
methodologies: MAP-Elites~\cite{mouret2015illuminating},
PAIR~\cite{chao2024jailbreaking}, and TAP~\cite{mehrotra2024tree} all
converge reliably, because at their core they are all
monotone-improvement operators acting on a bounded score function.

\begin{theorem}[Transferability]
\label{thm:transfer}
If $\|f-g\|_\infty \leq \delta$, then
$\{f > \tau+\delta\} \subseteq \{g > \tau\}$. If $\delta=0$, the basins
are identical.
\end{theorem}

The degree to which attacks transfer between models is governed by how
close the models' alignment deviation surfaces are in sup-norm. Models
with similar surfaces share overlapping vulnerability basins; distant
models do not.

\begin{theorem}[Authority Monotonicity]
\label{thm:authority}
If $f(y,\cdot)$ is monotone non-decreasing in authority for each fixed
indirection~$y$, then vulnerability is upward-closed, a critical threshold
$a_2^*$ exists, and the threshold curve is non-increasing in indirection.
\end{theorem}

\begin{theorem}[Gradient Ascent]
\label{thm:gradient}
If $f\colon E \to \R$ has nonzero Fr\'{e}chet derivative at~$x$ (in a
real inner product space~$E$), then there exist $v\in E$ and
$\varepsilon > 0$ with $f(x + \varepsilon v) > f(x)$.
\end{theorem}

At any non-critical point of the alignment deviation surface,
gradient-based attacks can always find a direction that increases the
score. The gradient serves as a reliable compass pointing toward higher
alignment deviation.


% ============================================================
\section{Cost Asymmetry}
\label{sec:cost}

Perhaps the most practically important consequence of the landscape
structure is the asymmetry it creates between attack and defense costs.

\begin{theorem}[Exponential Cost Asymmetry]
\label{thm:cost}
For grid resolution $N\geq 2$, attack step size $\delta > 0$, and
defense patch radius $0 < r < 1$:
\begin{enumerate}[nosep]
  \item Attack cost $\leq 1/\delta$ \textup{(dimension-independent)}.
  \item Defense cost $= N^d$ \textup{(exponential in dimension~$d$)}.
  \item The ratio $\delta \cdot N^d \to \infty$ as $d\to\infty$.
  \item For any $R > 0$, there exists $d_0$ such that the cost ratio
    exceeds $R$ for all $d \geq d_0$.
\end{enumerate}
\end{theorem}

\begin{proof}
The attack cost follows from \Cref{thm:convergence}: convergence takes at
most $\lceil 1/\delta\rceil$ steps, regardless of dimension. For the
defense cost, partitioning each axis into $N$ segments in $d$ dimensions
produces $N^d$ cells. By a pigeonhole argument (verified in Lean via
\texttt{Fintype.card\_le\_of\_surjective}), covering all cells requires
at least $N^d$ queries. The ratio $\delta \cdot N^d$ is monotone
increasing in~$d$ and diverges since $N^d \to \infty$ for $N \geq 2$.
\end{proof}

To get a sense of the numbers: at $d=2$ with $N=25$ and $\delta=0.01$,
the cost ratio is $6.25$---the defender pays roughly $6\times$ as much
as the attacker. At $d=10$, the ratio climbs to $\sim\!10^{12}$, making
comprehensive defense essentially intractable. The Lean artifact verifies
the $d=2$ calculation ($25^2 \times 0.01 = 6.25$) by \texttt{norm\_num}.


% ============================================================
\section{Experimental Validation}
\label{sec:experiments}

To test whether the theory describes real model behavior, we validate each
theorem against empirical findings from the Manifold of Failure
framework~\cite{bhatt2025manifold}, which maps the behavioral topology
of three LLMs using MAP-Elites over a 2D behavioral space (query
indirection $\times$ authority framing).

\subsection{Theory--Experiment Correspondence}

\begin{table}[t]
\centering
\small
\caption{Each theorem generates a testable prediction. In every case, the
empirical data confirms it.}
\label{tab:correspondence}
\begin{tabular}{@{}p{3.2cm}p{4.5cm}p{4.5cm}@{}}
\toprule
\textbf{Theorem} & \textbf{Predicts} & \textbf{Confirmed by} \\
\midrule
Basin Structure (\ref{thm:basin})
  & Basins have positive area
  & Heatmaps show extended colored regions, not isolated dots \\[4pt]
No Gap (\ref{thm:nogap})
  & Boundary is nonempty
  & Contour lines exist in all models with both safe and unsafe regions \\[4pt]
Perturbation Robustness (\ref{thm:lipschitz})
  & Adjacent cells have similar AD
  & Heatmaps are smooth; 3D surfaces are continuous \\[4pt]
Iterative Convergence (\ref{thm:convergence})
  & MAP-Elites/PAIR/TAP converge
  & Convergence curves are monotone and plateau \\[4pt]
Transferability (\ref{thm:transfer})
  & Close models share basins
  & Model spectrum: Llama ($0.93$) $\to$ GPT-OSS ($0.73$) $\to$ GPT-5-Mini ($0.47$) \\[4pt]
Authority Monotonicity (\ref{thm:authority})
  & Horizontal banding
  & Contour plots show bands at specific $a_2$ values \\[4pt]
Boundary Fixation (\ref{thm:main})
  & Defense can't clear boundary
  & GPT-5-Mini caps at $\tau=0.50$; no method breaches it \\[4pt]
Cost Asymmetry (\ref{thm:cost})
  & 2D tractable, high-$d$ intractable
  & 15K queries fill 63\% at $d\!=\!2$ \\
\bottomrule
\end{tabular}
\end{table}

\subsection{Model-Specific Predictions}

\paragraph{Llama-3-8B.} With a mean AD of $0.93$ and a basin rate of
$93.9\%$, this model's alignment deviation surface is a near-flat mesa.
In the language of \Cref{thm:lipschitz}, the Lipschitz constant~$L$ is
small, which means each vulnerability enjoys a large robustness radius.
The basin is essentially one massive connected region
(\Cref{thm:basin}), with only narrow safe corridors surviving at
specific authority levels. \Cref{thm:main} predicts that fixed boundary
points exist in these corridors, and because the corridors are so thin,
even small perturbations push prompts back onto the vulnerability plateau.

\paragraph{GPT-OSS-20B.} At a mean AD of $0.73$ and a basin rate of
$64.3\%$, this model presents a much more rugged landscape. The large
Lipschitz constant means small robustness radii
(\Cref{thm:lipschitz}), and the basin fragments into a complex mosaic
of safe and unsafe patches. \Cref{thm:authority} predicts horizontal
banding in this mosaic, and the contour plots confirm bands at
$a_2 \approx 0.25$--$0.35$ and $0.65$--$0.85$. \Cref{thm:main}
predicts fixed boundary points scattered throughout the safe--unsafe
mosaic, and the empirical ``bullseye'' patterns visible around
vulnerability peaks are exactly the contour structure one would expect
surrounding such points.

\paragraph{GPT-5-Mini.} This model reaches a peak AD of only $0.50$ with
a basin rate of $0\%$. If we set $\tau = 0.50$, the hypothesis
$U_\tau \neq \emptyset$ fails---there are no points strictly above
threshold---so \Cref{thm:main} does not apply. The theorem correctly
predicts that there is no impossibility here. The model's strategy is not
to eliminate the boundary but to make it harmless. At lower thresholds
like $\tau = 0.40$, \Cref{thm:basin} still applies (the model does
exhibit AD values around $0.47$ uniformly), but the ``basin'' at that
threshold is simply the entire space---uniformly moderate, never harmful.

\subsection{Quantitative Checks}

\paragraph{Robustness radius.} From the 3D surface plots, we estimate
$L \approx 5$ for Llama-3-8B (the surface changes by $\sim$0.5 over
$\sim$0.1 in behavioral coordinates). For a typical attack with
$\AD(p) = 0.93$ at $\tau = 0.50$, the predicted robustness radius is
$r(p) = (0.93-0.50)/5 = 0.086$. On a $25\times 25$ grid with cell width
$0.04$, this covers roughly 2 cells in each direction---consistent with
what we see in the heatmaps, where adjacent cells consistently take
similar values.

\paragraph{Convergence bound.} The typical per-step improvement observed
in the convergence curves is $\delta \approx 0.001$--$0.01$, which
predicts convergence in $100$--$1000$ steps. The empirical curves plateau
within $\sim$5000 iterations, which is consistent with the upper bound
once you account for the fact that early iterations achieve larger gains
than the steady-state $\delta$.

\paragraph{Cost ratio.} At $d=2$, $N=25$, $\delta = 0.01$, the
theoretical ratio is $0.01 \times 625 = 6.25$. In practice, the
framework uses 15,000 queries to fill 394 out of 625 cells ($63\%$
coverage), averaging $\sim$38 queries per cell. The attacker needs
$\sim$100 queries to converge ($1/\delta$). The empirical ratio is
therefore $394/100 \approx 4$, which is in the same ballpark as the
theoretical $6.25$.


% ============================================================
\section{The Engineering Prescription}
\label{sec:engineering}

\Cref{thm:main} tells us the boundary cannot be eliminated. That still
leaves the question of what to do about it. The theory points toward three
concrete strategies.

\paragraph{Strategy 1: Make the boundary shallow.}
If you set $\tau$ high enough that $f(z) = \tau$ corresponds to benign
behavior, the theorem still holds but ceases to matter. GPT-5-Mini
exemplifies this approach: its boundary points have $\AD = 0.50$, which
corresponds to moderate-level refusals rather than harmful compliance.
The boundary exists; it just is not dangerous.

\paragraph{Strategy 2: Reduce the Lipschitz constant.}
\Cref{thm:lipschitz} tells us the robustness radius scales as
$(f(p)-\tau)/L$. A smoother alignment surface (smaller~$L$) means
that any remaining vulnerabilities extend over wider regions---but it also
means the surface is more predictable and easier to monitor. The contour
plots for GPT-5-Mini illustrate this well: its surface barely varies,
with $\AD$ confined to the range $[0.39, 0.50]$.

\paragraph{Strategy 3: Reduce the effective dimension.}
\Cref{thm:cost} shows that defense cost grows as $N^d$. By constraining
the prompt interface---through standardized formats, restricted API
parameters, or bounded context lengths---one effectively reduces~$d$,
making the behavioral space more tractable to monitor and patrol.


% ============================================================
\section{The Lean Artifact}
\label{sec:artifact}

\subsection{Verification Details}

The complete theory has been machine-verified in Lean~4.28.0 using
Mathlib~v4.28.0. The artifact is organized into:
\begin{itemize}[nosep]
  \item 10 core theory files (MoF\_01--MoF\_10)
  \item 10 cost theory files (MoF\_Cost\_01--MoF\_Cost\_10)
  \item 10 advanced theory files (MoF\_Adv\_01--MoF\_Adv\_10)
  \item 3 capstone files (MasterTheorem, Euclidean instantiation, verification)
\end{itemize}

\noindent Every theorem in the artifact depends on exactly three axioms:
\begin{itemize}[nosep]
  \item \texttt{propext}: propositional extensionality ($P \leftrightarrow Q \implies P = Q$)
  \item \texttt{Classical.choice}: the axiom of choice
  \item \texttt{Quot.sound}: quotient soundness
\end{itemize}

\noindent These are Lean's standard foundational axioms, shared by all of
Mathlib. No additional axioms are introduced.

\subsection{Proof Map}

The master theorem's proof chain draws on six Mathlib lemmas:

\begin{center}
\small
\begin{tabular}{@{}ll@{}}
\toprule
\textbf{Lean name} & \textbf{Says} \\
\midrule
\texttt{isClosed\_diagonal} & Diagonal is closed in Hausdorff spaces \\
\texttt{IsClosed.preimage} & Preimage of closed under continuous is closed \\
\texttt{IsClosed.closure\_subset\_iff} & Closed set containing $S$ contains $\cl{S}$ \\
\texttt{isClopen\_iff} & Clopen sets in connected spaces are $\emptyset$ or $X$ \\
\texttt{closure\_minimal} & $\cl{S} \subseteq T$ when $S\subseteq T$ and $T$ closed \\
\texttt{isClosed\_le} & $\{f \leq g\}$ is closed for continuous $f,g$ \\
\bottomrule
\end{tabular}
\end{center}

\subsection{Generality}

All theorems are stated over arbitrary topological spaces
(\texttt{\{X~:~Type*\} [TopologicalSpace~X]}) unless they specifically
require metric or measure structure. A separate Euclidean instantiation
file verifies that $\mathbb{R}^n$
(\texttt{EuclideanSpace~$\mathbb{R}$~(Fin~d)}) satisfies all required
typeclasses for every dimension~$d$.


% ============================================================
\section{Limitations}
\label{sec:limitations}

\paragraph{The continuous model is an approximation.} Real prompts are
discrete token sequences, and the behavioral space is finite. Our theory
operates on a continuous relaxation of this discrete reality. The
approximation is empirically supported---Gaussian process fits to the
behavioral data achieve $R^2 > 0.85$~\cite{bhatt2025manifold}---but it
is not formally guaranteed to hold.

\paragraph{The impossibility is at the boundary.} The fixed points
produced by the theorem satisfy $f(z) = \tau$ exactly, which is the
mildest possible kind of failure. The result constrains perfection, not
adequacy. A defense that reduces alignment deviation to
$\tau + \varepsilon$ everywhere except at a handful of boundary points is
entirely consistent with the theorem.

\paragraph{The cost asymmetry assumes brute-force search.}
\Cref{thm:cost} counts grid cells and derives a lower bound based on
exhaustive enumeration. Defenders who use learning-based approaches that
generalize across the space may be able to sidestep the exponential
lower bound by exploiting structure that the theorem does not account for.

\paragraph{Single-turn only.} The manifold as we model it is static, but
in practice it may shift across conversation turns as context accumulates.
Extending the theory to time-varying or history-dependent manifolds
remains an open direction for future work.

\paragraph{The theory does not explain model-specific topology.} Our
theorems describe what must be true of the landscape given continuity and
connectedness, but they do not explain \emph{why} Llama-3-8B has a
near-flat vulnerability plateau while GPT-5-Mini does not. Those
differences stem from architecture, training data, and alignment
procedure---empirical questions that lie outside the scope of our
mathematical framework.


% ============================================================
\section{Conclusion}
\label{sec:conclusion}

We have shown that any continuous, utility-preserving defense must leave
boundary points of the alignment deviation surface fixed in place. The
proof is five steps long, with each step drawing on exactly one
hypothesis. The three hypotheses---continuity, utility preservation, and
connectedness---are all individually necessary, as demonstrated by
counterexamples for each. Together, they establish a defense trilemma:
no defense can simultaneously be continuous, utility-preserving, and
complete.

Nine supporting theorems flesh out the landscape that this impossibility
constrains. Vulnerability basins are open sets with positive measure.
The boundary between safe and unsafe is unavoidable. Successful attacks
are robust to perturbation and converge reliably under iteration. Transfer
between models is controlled by the sup-norm distance between their
alignment deviation surfaces. And the cost of comprehensive defense grows
exponentially in the dimension of the behavioral space, while attack cost
does not. Each of these results finds confirmation in empirical data from
three frontier language models.

The practical upshot is not that defense is futile, but that it requires
calibration. Make the boundary shallow by raising~$\tau$ so that
boundary-level behavior is benign. Make the surface smooth by
reducing~$L$ so that the landscape is predictable. Make the space
low-dimensional by constraining the prompt interface so that it is
tractable to patrol. GPT-5-Mini's hard ceiling at $\AD = 0.50$ shows
what this prescription looks like in practice.

The full theory is machine-verified in Lean~4: 33 files, approximately
200 theorems, zero \texttt{sorry} statements, and three standard axioms.
The artifact is open-sourced for the community to build on.


% ============================================================
\bibliographystyle{plain}
\begin{thebibliography}{99}

\bibitem{alon2023detecting}
G.~Alon and M.~Kamfonas.
\newblock Detecting language model attacks with perplexity.
\newblock \emph{arXiv preprint arXiv:2308.14132}, 2023.

\bibitem{bagnall2019certifying}
A.~Bagnall and G.~Stewart.
\newblock Certifying the true error: Machine learning in {Coq} with
  verified generalization guarantees.
\newblock \emph{Proceedings of the AAAI Conference on Artificial Intelligence}, 2019.

\bibitem{bai2022constitutional}
Y.~Bai et~al.
\newblock Constitutional {AI}: Harmlessness from {AI} feedback.
\newblock \emph{arXiv preprint arXiv:2212.08073}, 2022.

\bibitem{bhatt2025manifold}
M.~Bhatt, S.~Munshi, V.~S.~Narajala, I.~Habler, A.~Al-Kahfah, K.~Huang, and B.~Gatto.
\newblock Manifold of failure: Behavioral attraction basins in language
  models.
\newblock \emph{arXiv preprint arXiv:2602.22291v2}, 2026.

\bibitem{bolte2024}
J.~Bolte et~al.
\newblock Formal verification of optimization algorithms.
\newblock \emph{Mathematical Programming}, 2024.

\bibitem{chao2024jailbreaking}
P.~Chao, A.~Robey, E.~Dobriban, H.~Hassani, G.~J. Pappas, and E.~Wong.
\newblock Jailbreaking black box large language models in twenty queries.
\newblock \emph{arXiv preprint arXiv:2310.08419}, 2024.

\bibitem{cohen2019certified}
J.~Cohen, E.~Rosenfeld, and J.~Z. Kolter.
\newblock Certified adversarial robustness via randomized smoothing.
\newblock \emph{Proceedings of ICML}, 2019.

\bibitem{goodfellow2014explaining}
I.~J. Goodfellow, J.~Shlens, and C.~Szegedy.
\newblock Explaining and harnessing adversarial examples.
\newblock \emph{arXiv preprint arXiv:1412.6572}, 2014.

\bibitem{inan2023llama}
H.~Inan et~al.
\newblock {Llama Guard}: {LLM}-based input-output safeguard for
  human-{AI} conversations.
\newblock \emph{arXiv preprint arXiv:2312.06674}, 2023.

\bibitem{katz2017reluplex}
G.~Katz, C.~Barrett, D.~L. Dill, K.~Julian, and M.~J. Kochenderfer.
\newblock Reluplex: An efficient {SMT} solver for verifying deep neural
  networks.
\newblock \emph{Proceedings of CAV}, 2017.

\bibitem{mehrotra2024tree}
A.~Mehrotra et~al.
\newblock Tree of attacks: Jailbreaking black-box {LLMs} automatically.
\newblock \emph{arXiv preprint arXiv:2312.02119}, 2024.

\bibitem{mouret2015illuminating}
J.-B. Mouret and J.~Clune.
\newblock Illuminating search spaces by mapping elites.
\newblock \emph{arXiv preprint arXiv:1504.04909}, 2015.

\bibitem{samvelyan2024rainbow}
M.~Samvelyan et~al.
\newblock Rainbow teaming: Open-ended generation of diverse adversarial
  prompts.
\newblock \emph{Advances in Neural Information Processing Systems}, 37, 2024.

\bibitem{singh2019abstract}
G.~Singh, T.~Gehr, M.~P{\"u}schel, and M.~Vechev.
\newblock An abstract domain for certifying neural networks.
\newblock \emph{Proceedings of POPL}, 2019.

\bibitem{szegedy2013intriguing}
C.~Szegedy et~al.
\newblock Intriguing properties of neural networks.
\newblock \emph{arXiv preprint arXiv:1312.6199}, 2013.

\bibitem{tsipras2018robustness}
D.~Tsipras, S.~Santurkar, L.~Engstrom, A.~Turner, and A.~Madry.
\newblock Robustness may be at odds with accuracy.
\newblock \emph{arXiv preprint arXiv:1805.12152}, 2018.

\bibitem{wolpert1997no}
D.~H. Wolpert and W.~G. Macready.
\newblock No free lunch theorems for optimization.
\newblock \emph{IEEE Transactions on Evolutionary Computation}, 1(1):67--82, 1997.

\bibitem{zou2023universal}
A.~Zou, Z.~Wang, N.~Carlini, M.~Nasr, J.~Z. Kolter, and M.~Fredrikson.
\newblock Universal and transferable adversarial attacks on aligned
  language models.
\newblock \emph{arXiv preprint arXiv:2307.15043}, 2023.

\end{thebibliography}

\end{document}
